% =========================================================================
% ปฏิบัติการที่ 4: การสร้างปฏิสัมพันธ์กับวัตถุ 3 มิติในพื้นที่เสมือน
% =========================================================================

\chapter{การสร้างปฏิสัมพันธ์กับวัตถุ 3 มิติในพื้นที่เสมือน}
\label{ch:lab04}

\noindent{\large\bfseries\color{cyberblue} การประยุกต์ Raycasting เพื่อการตรวจจับการชน การยกย้าย และการหมุนโมเดล}

\vspace{0.4cm}

\begin{labobjectivebox}{การสร้างปฏิสัมพันธ์กับวัตถุ 3 มิติในพื้นที่เสมือน}
\textbf{วัตถุประสงค์การเรียนรู้}
\begin{enumerate}[leftmargin=1.8cm, itemsep=2pt]
    \item เข้าใจทฤษฎีการฉายรังสีตรวจจับการชน (Raycasting Collision Detection)
    \item สามารถแปลงพิกัด Normalized Device Coordinates (NDC) จากมือสู่ระนาบ 3 มิติ
    \item สามารถเขียนโปรแกรมจับ ยก ย้าย และปล่อยวัตถุในฉากเสมือนจริงได้อย่างราบรื่น
\end{enumerate}

\vspace{4pt}
\textbf{เครื่องมือและอุปกรณ์จำเป็น} \enspace Three.js Raycaster, เวกเตอร์พิกัดมือจากการทดลองที่ 3
\end{labobjectivebox}

\section{หลักการและทฤษฎีพื้นฐาน}

การสร้างปฏิสัมพันธ์ระหว่างผู้ใช้กับวัตถุในสภาพแวดล้อม 3 มิติโดยไม่มีอุปกรณ์สัมผัสกายภาพ อาศัยกระบวนการฉายรังสีเสมือน (Raycasting) ซึ่งเป็นการสร้างรังสีเรขาคณิต (Mathematical Ray) ที่มีจุดกำเนิด $\mathbf{O}$ และเวกเตอร์ทิศทาง $\mathbf{D}$:
\begin{equation}
\mathbf{R}(t) = \mathbf{O} + t\mathbf{D}, \quad t \ge 0
\end{equation}
ในระบบความจริงเสริม พิกัดจากกล้องเว็บแคมจะถูกแปลงเข้าสู่ระบบพิกัดอุปกรณ์มาตรฐาน (Normalized Device Coordinates: NDC) โดยมีขอบเขต $[-1.0, 1.0]$ ทั้งแกน $X$ และ $Y$:
\begin{equation}
x_{\text{ndc}} = 2 \cdot x_{\text{hand}} - 1, \quad y_{\text{ndc}} = -(2 \cdot y_{\text{hand}} - 1)
\end{equation}
เมื่อนำเวกเตอร์ $(x_{\text{ndc}}, y_{\text{ndc}})$ ส่งผ่านกล้องเสมือน (Camera Projection Matrix) ตัวเรนเดอร์จะยิงรังสีทะลุเข้าไปในฉากทัศน์ 3 มิติ และคำนวณจุดตัด (Intersection) กับรูปทรงเรขาคณิตของวัตถุ หากเกิดจุดตัดและผู้ใช้ทำท่าทางจีบนิ้ว (Pinch) วัตถุจะถูกผูกโยงพิกัดเข้ากับมือของผู้ใช้ ทำให้สามารถเคลื่อนย้ายวัตถุไปตามทิศทางการเคลื่อนไหวของมือได้อย่างสมจริง

\begin{conceptbox}{สาระสำคัญของปฏิบัติการที่ 4}
การเข้าใจโครงสร้างทางคณิตศาสตร์และการประมวลผลเชิงพื้นที่ ช่วยให้นักศึกษาสามารถออกแบบระบบเสมือนจริงที่ตอบสนองต่อผู้ใช้งานได้อย่างเป็นธรรมชาติและแม่นยำสูง ทั้งยังสามารถต่อยอดสู่การพัฒนาแอปพลิเคชันทางวิทยาศาสตร์และอุตสาหกรรมได้อย่างมีประสิทธิภาพ
\end{conceptbox}

\section{ขั้นตอนการปฏิบัติการและการทดลอง}

ให้นักศึกษาดำเนินการสร้างไฟล์โครงการและเขียนคำสั่งตามลำดับขั้นตอนต่อไปนี้ โดยตรวจสอบความถูกต้องของไลบรารีที่เรียกใช้งานและเส้นทางไฟล์

\begin{codebox}{โครงสร้างโค้ดหลักของปฏิบัติการที่ 4}
\begin{lstlisting}[language=HTML]
// การสร้าง Raycaster และแปลงพิกัดมือเป็น NDC
const raycaster = new THREE.Raycaster();
const mouseNDC = new THREE.Vector2();
let selectedObject = null;

function updateHandInteraction(cursor, isPinching) {
    // แปลงพิกัดมือ [0, 1] เป็นพิกัด NDC [-1, 1]
    mouseNDC.x = (cursor.x * 2) - 1;
    mouseNDC.y = -(cursor.y * 2) + 1;

    // ยิงรังสีจากตำแหน่งกล้องผ่านพิกัด NDC
    raycaster.setFromCamera(mouseNDC, camera);
    const intersects = raycaster.intersectObjects(scene.children);

    if (intersects.length > 0) {
        const hit = intersects[0].object;
        if (isPinching && !selectedObject) {
            selectedObject = hit;
            selectedObject.material.emissive.setHex(0x06b6d4); // เรืองแสงเมื่อถูกจับ
        }
    }

    // หากมีการจับวัตถุอยู่ ให้เคลื่อนย้ายวัตถุตามพิกัด 3D ของรังสี
    if (selectedObject) {
        if (isPinching) {
            const targetPos = raycaster.ray.at(3.0, new THREE.Vector3());
            selectedObject.position.copy(targetPos);
        } else {
            // ปล่อยวัตถุเมื่อเลิกจีบนิ้ว
            selectedObject.material.emissive.setHex(0x000000);
            selectedObject = null;
        }
    }
}
\end{lstlisting}
\end{codebox}

\section{การทดสอบระบบและการบันทึกผล}

เมื่อรันโปรแกรมบนเซิร์ฟเวอร์จำลอง (Local Development Server) ให้ทำการทดสอบการทำงานตามเกณฑ์ประเมินในตารางต่อไปนี้ พร้อมบันทึกผลการสังเกตการณ์ลงในรายงานผลการทดลอง

\begin{table}[htbp]
\centering
\small
\begin{tabularx}{\textwidth}{c X c Y}
\toprule
\textbf{ลำดับ} & \textbf{รายการทดสอบและพฤติกรรมที่คาดหวัง} & \textbf{เกณฑ์ผ่าน} & \textbf{ผลการทดสอบ} \\
\midrule
1 & การเริ่มต้นระบบและการเข้าถึงสิทธิ์ฮาร์ดแวร์ & ภายใน 3 วินาที & ผ่าน / ไม่ผ่าน \\
2 & ความเสถียรของการตรวจจับและการคำนวณพิกัด & อัตราความแม่นยำ $> 90\%$ & ผ่าน / ไม่ผ่าน \\
3 & อัตราการแสดงผลกราฟิกต่อเนื่อง (Framerate) & ไม่ต่ำกว่า 55 FPS & ผ่าน / ไม่ผ่าน \\
4 & การตอบสนองต่อปฏิสัมพันธ์เชิงพื้นที่ & ความหน่วง $< 40$ ms & ผ่าน / ไม่ผ่าน \\
\bottomrule
\end{tabularx}
\caption{ตารางประเมินผลการทำงานของระบบในปฏิบัติการที่ 4}
\label{tab:eval_lab04}
\end{table}

\section{ภารกิจท้าทายและการต่อยอด}

\begin{activitybox}{กิจกรรมส่งเสริมทักษะขั้นสูง}
ให้นักศึกษาเพิ่มคุณสมบัติการหมุนวัตถุ (Rotation) ตามการเอียงของแนวระนาบฝ่ามือ โดยใช้เวกเตอร์ระหว่างจุด Landmark 0 (ข้อมือ) และ Landmark 9 (โคนนิ้วกลาง)
\end{activitybox}

\section{คำถามท้ายการทดลอง}

ให้นักศึกษาตอบคำถามเชิงวิเคราะห์ต่อไปนี้ลงในสมุดบันทึกปฏิบัติการ

\begin{enumerate}[leftmargin=1.8cm, itemsep=6pt]
    \item เพราะเหตุใดสูตรการแปลงพิกัด $y_{\text{hand}}$ สู่ $y_{\text{ndc}}$ จึงต้องติดเครื่องหมายลบด้านหน้า
    \item การคำนวณการชนด้วย Raycasting กับโมเดลที่มีโพลีกอนจำนวนมาก (High-poly models) ส่งผลต่อประสิทธิภาพอย่างไร และมีวิธีแก้ไขอย่างไร
    \item จงอธิบายหลักการของ Bounding Box (AABB) ในการเพิ่มความเร็วให้แก่อัลกอริทึมตรวจจับการชน
\end{enumerate}

